\documentclass[11pt,a4paper]{article}
\usepackage[margin=23mm]{geometry}
\usepackage{iftex}
\ifPDFTeX
  \usepackage[T2A]{fontenc}
  \usepackage[utf8]{inputenc}
\else
  \usepackage{fontspec}
  \IfFileExists{cmunrm.otf}{%
    \setmainfont{cmunrm.otf}[BoldFont=cmunbx.otf,
      ItalicFont=cmunti.otf,BoldItalicFont=cmunbi.otf]
  }{%
    \setmainfont{Times New Roman}
  }
\fi
\usepackage[english,russian]{babel}
\usepackage{amsmath,amssymb}
\usepackage[unicode,hidelinks]{hyperref}
\setlength{\parindent}{0pt}
\setlength{\parskip}{4pt}
\setlength{\emergencystretch}{3em}
\begin{document}
\begin{center}
{\large\bfseries Топологический анализ спектра\\
пространственно-временных рядов}
\end{center}

\section*{1. Данные и задача}
Даны $q$ рядов, измеренных датчиками в точках $s_i$:
\[
 Y_{in}=x(s_i,t_n),\qquad t_n=n\Delta t,\quad
 1\leq i\leq q,\quad 0\leq n<N.
\]

\section*{2. SSA: разложение одного ряда}
SSA состоит из построения матрицы окон,
ее сингулярного разложения и восстановления выбранных компонент.
Для ряда $y_0,\ldots,y_{N-1}$ выбираем длину окна $1<L<N$,
полагаем $K=N-L+1$ и записываем
\[
 H=\begin{pmatrix}
 y_0&y_1&\cdots&y_{K-1}\\
 y_1&y_2&\cdots&y_K\\
 \vdots&\vdots&&\vdots\\
 y_{L-1}&y_L&\cdots&y_{N-1}
 \end{pmatrix}\in\mathbb R^{L\times K}.
\]
Столбец $b$ есть окно $(y_b,\ldots,y_{b+L-1})^T$.

Для любой вещественной матрицы ранга $R$ существует SVD:
\[
 H=U\Sigma V^T=\sum_{r=1}^{R}H_r,\qquad
 H_r=\sigma_ru_rv_r^T,\quad \sigma_1\geq\cdots\geq\sigma_R>0.
\]
$u_r\in\mathbb R^L$, $v_r\in\mathbb R^K$ --- единичные векторы;
наборы $u_r$ и $v_r$ ортонормированы. $\sigma_r$ --- сингулярные числа.

\textbf{Что выбирается при восстановлении?}
$G\subseteq\{1,\ldots,R\}$ --- множество номеров сохраняемых слагаемых:
\[
 \widetilde H=\sum_{r\in G}H_r.
\]
Например, $G=\{1,\ldots,d\}$ сохраняет первые $d$ слагаемых;
$G=\{r_1,r_2\}$ сохраняет выбранную пару. SVD не определяет,
какая группа соответствует конкретному колебанию.

В исходной матрице все элементы с $a+b=n$ равны $y_n$.
После отбора слагаемых они могут различаться, поэтому берем их среднее:
\[
 \widetilde y_n=\frac1{w_n}\sum_{a+b=n}\widetilde H_{ab}.
\]
Здесь $w_n$ --- число пар $(a,b)$ в пределах матрицы с $a+b=n$.
Это диагональное усреднение, завершающее SSA.

\section*{3. MSSA: совместное разложение нескольких рядов}
MSSA (Multichannel Singular Spectrum Analysis) применяет одно SVD
к матрицам окон всех $q$ рядов. Алгоритм описан в
\cite[раздел 1.1, с.~294--295]{mssa}.
Для каждого ряда $Y_{i0},\ldots,Y_{i,N-1}$ строим
\[
 (H_i)_{ab}=Y_{i,a+b},\qquad
 0\leq a<L,\quad 0\leq b<K,\quad K=N-L+1.
\]
В горизонтальной MSSA эти матрицы ставятся рядом:
\[
 \mathcal H=[H_1\ \cdots\ H_q]
 =\left(\begin{array}{cccc|c|cccc}
 Y_{1,0}&Y_{1,1}&\cdots&Y_{1,K-1}&\cdots&
 Y_{q,0}&Y_{q,1}&\cdots&Y_{q,K-1}\\
 Y_{1,1}&Y_{1,2}&\cdots&Y_{1,K}&\cdots&
 Y_{q,1}&Y_{q,2}&\cdots&Y_{q,K}\\
 \vdots&\vdots&&\vdots&&\vdots&\vdots&&\vdots\\
 Y_{1,L-1}&Y_{1,L}&\cdots&Y_{1,N-1}&\cdots&
 Y_{q,L-1}&Y_{q,L}&\cdots&Y_{q,N-1}
 \end{array}\right)
 \in\mathbb R^{L\times qK}.
\]
Каждый блок содержит $K$ окон одного датчика; строки задают позиции
отсчетов внутри окна. SVD вычисляется для всей $\mathcal H$:
\[
 \mathcal H=U\Sigma V^T,\qquad
 \mathcal H\mathcal H^T=\sum_{i=1}^qH_iH_i^T.
\]
Левые сингулярные векторы $u_r\in\mathbb R^L$ выбираются совместно
по всем рядам. Для выбранных номеров $G$ получаем
\begin{gather*}
 P_G=\sum_{r\in G}u_ru_r^T,\\
 \widetilde{\mathcal H}=P_G\mathcal H
 =[P_GH_1\ \cdots\ P_GH_q],\qquad
 \widetilde H_i=P_GH_i.
\end{gather*}
В каждом блоке отдельно усредняем элементы с $a+b=n$:
\[
 \widetilde Y_{in}
 =\frac1{w_n}\sum_{\substack{a+b=n\\0\leq a<L,\;0\leq b<K}}
 (\widetilde H_i)_{ab},\qquad 0\leq n<N.
\]
Получаем $q$ восстановленных рядов с одним общим проектором $P_G$.
При отдельных SSA каждый ряд имел бы свой проектор.
Координаты датчиков $s_i$ в построение $\mathcal H$ не входят.

\section*{4. Ранг гармоники}
Рассмотрим гармонику с амплитудой $A$ и начальной фазой $\alpha$:
\[
 y_n=A\cos(\phi n+\alpha),\qquad A\ne0,\quad \phi=\omega\Delta t.
\]
По формуле косинуса суммы
\begin{align*}
 y_{a+b}
 &=A\cos\bigl(\phi a+(\phi b+\alpha)\bigr)\\
 &=A\cos(\phi b+\alpha)\cos(\phi a)
   -A\sin(\phi b+\alpha)\sin(\phi a).
\end{align*}
Определим два фиксированных вектора длины $L$:
\begin{align*}
 c&=(1,\cos\phi,\ldots,\cos((L-1)\phi))^T,\\
 s&=(0,\sin\phi,\ldots,\sin((L-1)\phi))^T.
\end{align*}
Тогда каждый столбец матрицы окон имеет вид
\[
 H_{:,b}=A\cos(\phi b+\alpha)c-A\sin(\phi b+\alpha)s.
\]
Все столбцы лежат в пространстве $\operatorname{span}\{c,s\}$,
поэтому $\operatorname{rank}H\leq2$.
При $L,K\geq2$ первые две строки и два столбца дают определитель
\[
 A^2\det\begin{pmatrix}
 \cos\alpha&\cos(\alpha+\phi)\\
 \cos(\alpha+\phi)&\cos(\alpha+2\phi)
 \end{pmatrix}=-A^2\sin^2\phi.
\]
Если $0<\phi<\pi$, он ненулевой, значит ранг не меньше двух.
Вместе с предыдущей оценкой получаем $\operatorname{rank}H=2$.

\textbf{Одна гармоника занимает два направления SVD.}
Одно слагаемое ранга один обычно не восстанавливает ее целиком.
Для суммы $h$ гармоник все столбцы лежат в пространстве $h$ косинусных
и $h$ синусных векторов, поэтому ранг не больше $2h$.

При нескольких гармониках SVD может смешивать разные частоты в одной
компоненте. Поэтому пара компонент SVD не всегда соответствует
одной гармонике.

\section*{5. Пространственные фазы и траектория измерений}
Для одной частоты два датчика дают
\begin{align*}
 x_1&=a_1\cos\theta,\qquad \theta=\omega t,\quad a_1,a_2>0,\\
 x_2&=a_2\cos(\theta+\delta)
     =a_2\cos\delta\cos\theta-a_2\sin\delta\sin\theta.
\end{align*}
$\delta$ --- фазовый сдвиг между датчиками; начало времени выбрано
по первому датчику. Объединяем показания:
\[
 X=\begin{pmatrix}x_1\\x_2\end{pmatrix}
 =\underbrace{\begin{pmatrix}a_1&0\\a_2\cos\delta&-a_2\sin\delta\end{pmatrix}}_B
 \begin{pmatrix}\cos\theta\\\sin\theta\end{pmatrix},\qquad
 \det B=-a_1a_2\sin\delta.
\]
При $\sin\delta\ne0$ матрица обратима: окружность переходит в эллипс,
а фаза определяется по $X$ однозначно по модулю $2\pi$.
При $\sin\delta=0$ получается отрезок, и $\theta,-\theta$
дают одинаковые показания.

Для нескольких частот принимаем модель с постоянными амплитудами и фазами:
\begin{align*}
 x_i(t)&=\sum_{j=1}^h a_{ij}\cos(\omega_jt+\varphi_{ij})\\
 &=\sum_{j=1}^h\left[
 a_{ij}\cos\varphi_{ij}\cos\theta_j
 -a_{ij}\sin\varphi_{ij}\sin\theta_j\right],\qquad \theta_j=\omega_jt.
\end{align*}
Коэффициенты при косинусе и синусе частоты $j$ образуют столбцы
$2j-1$ и $2j$ матрицы $B$:
\begin{gather*}
 p=(\cos\theta_1,\sin\theta_1,\ldots,\cos\theta_h,\sin\theta_h)^T,\\
 B_{i,2j-1}=a_{ij}\cos\varphi_{ij},\qquad
 B_{i,2j}=-a_{ij}\sin\varphi_{ij},\\
 X=Bp.
\end{gather*}
Каждая пара координат $p$ лежит на окружности. Все возможные фазы
образуют произведение $h$ окружностей, тор $\mathbb T^h$.
Условие $\operatorname{rank}B=2h$ позволяет восстановить $p$ по $X$
и сохраняет различимость всех точек этого тора.

При $\omega_2=2\omega_1$ имеем $\theta_2=2\theta_1\pmod{2\pi}$:
свободная фаза одна, и фазовое замыкание есть окружность.
При иррациональном $\omega_2/\omega_1$ замыкание непрерывной траектории
есть $\mathbb T^2$ \cite[следствие 2.3]{torus}.
Число спектральных линий может превышать число независимых фаз.
Дискретизация и конечная запись могут ограничивать покрытие фаз.

Задержка $\tau$ заменяет фазовый сдвиг $\varphi_{ij}$ на
$\varphi_{ij}-\omega_j\tau$ и добавляет строку в матрицу наблюдения:
\[
 x_i(t-\tau)=\sum_j a_{ij}\left[
 \cos(\varphi_{ij}-\omega_j\tau)\cos\theta_j
 -\sin(\varphi_{ij}-\omega_j\tau)\sin\theta_j\right].
\]
Для одной частоты пара $(x_1(t),x_1(t-\tau))$ сохраняет фазу
при $\sin(\omega\tau)\ne0$.

\section*{6. Ошибка ряда и ошибка матрицы}
Для ошибки $e_n=y_n-\widetilde y_n$ унитарное преобразование Фурье дает
\[
 \widehat e_k=\frac1{\sqrt N}\sum_{n=0}^{N-1}e_ne^{-2\pi ikn/N},\qquad
 \sum_{n=0}^{N-1}|e_n|^2=\sum_{k=0}^{N-1}|\widehat e_k|^2.
\]
Это равенство Парсеваля: сумма квадратов сохраняется при смене
ортонормированного базиса. Оно действует для полного комплексного
спектра. Для нескольких датчиков равенства складываются.

В матрице SSA отсчет $y_n$ повторяется $w_n$ раз, поэтому
\[
 \|H\|_F^2=\sum_nw_ny_n^2=\sum_{r=1}^R\sigma_r^2,\qquad
 \|H-\widetilde H\|_F^2=\sum_{r\notin G}\sigma_r^2.
\]
$\|M\|_F^2$ обозначает сумму квадратов элементов матрицы $M$.
Ошибка матрицы SSA и ошибка восстановленного ряда различаются;
последнюю считаем после диагонального усреднения.

\paragraph{Граница выводов.}
SSA/MSSA определены для произвольных данных. Выводы о фазовых окружностях
и торах относятся к точной сумме стационарных гармоник.
Оценка ошибки разделения частот при шуме здесь не выведена,
поэтому устойчивость восстановленной фазовой геометрии пока не обоснована.

\begin{thebibliography}{2}
\bibitem{mssa}
N.~Golyandina, D.~Stepanov.
\emph{SSA-based approaches to analysis and forecast of multidimensional
 time series}. Proceedings of the 5th St. Petersburg Workshop on
Simulation, 2005, с.~293--298; алгоритм: раздел 1.1, с.~294--295.
\url{https://www.gistatgroup.com/gus/mssa2.pdf}.
\bibitem{torus}
H.~Gakhar, J.~A.~Perea.
\emph{Sliding window persistence of quasiperiodic functions}.
arXiv:2103.04540v3, раздел 2.1.
\url{https://arxiv.org/html/2103.04540v3}.
\end{thebibliography}
\end{document}
